Artificial intelligence has become an infrastructure problem. Training and
serving frontier models at scale is no longer limited only by the accelerator at
the center of the rack, but by how many accelerators can be wired together
efficiently, reliably, and within a fixed power envelope. That wiring is
increasingly optical, and the optics (especially the lasers inside them) have
become a first-order lever on the cost, power, and reliability of the whole
system.

This book is a concise technical overview of that layer: the short-reach optical
interconnects that stitch together AI datacenters, from in-package optical I/O
out to intra-rack links, deliberately setting aside the 2--10~km campus links that
belong to coherent optics (\Cref{sec:reach}). It concentrates on high-speed
signaling (including multilevel PAM4/PAM8 design and channel equalization),
optical devices, wavelength control, packaging and fabric context, and a compact
productization path from requirements to controlled ramp. Failure analysis
closes the loop (\Cref{ch:failure-modes}). Detailed qualification, manufacturing,
and operations catalogs live in the appendices. Fabric survey context lives in
\Cref{app:ai-fabric-context}.

\newthought{How to use this handbook.} The book supports three modes. Pick one
and stay in it until you switch deliberately.

\begin{description}
\item[Deep learning] Read the body chapters from requirements downward
      (\Cref{ch:role} through \Cref{ch:productization}, including multilevel
      signaling in \Cref{ch:multilevel}, equalization in
      \Cref{ch:equalization}, packaging in \Cref{ch:packaging}, and HPC rack
      architecture in \Cref{ch:hpc}), then the failure-analysis method
      (\Cref{ch:failure-modes}). Every major chapter asks: How does it work?
      How is it measured, and what uncertainty does the measurement remove?
      How does it fail? How is it debugged?
\item[Interview preparation] Follow the one-week plan
      (\Cref{sec:interview-week-plan}), review the printable chapter cheat sheets
      (\Cref{app:interview-cheatsheets}), practice cases
      (\Cref{app:case-studies}), drill thirty-second answers and tradeoff
      questions (\Cref{app:interview-frameworks}), review the Top~25 index
      (\Cref{app:top25}), and read how Staff judgment works
      (\Cref{app:staff-thinking}). Prioritize signaling, devices, packaging, and
      fabric architecture over readiness terminology.
\item[Incident / problem solving] Open the wall-chart trees
      (\Cref{app:decision-trees}), failure handbook
      (\Cref{ch:failure-modes}), case (\Cref{app:case-studies}),
      Staff judgment (\Cref{app:staff-thinking}), or the productization spine
      (\Cref{ch:productization,tab:ladder}) and its appendix depth
      (\Cref{app:reliability-reference,app:manufacturing-reference,app:ai-fabric-context}).
\end{description}

Start from the requirement, not the component. Move downward from system
requirements to architecture, subsystem, component, and only then to the
physics needed to make the decision. A choice at one layer constrains the next.
The aim is operational judgment: reduce uncertainty enough to make the next
decision, then name the control that prevents recurrence.

\newthought{On sources.} The industry moves quickly. Where the text cites specific
products or figures (co-packaged-optics programs, per-lane roadmaps, energy-per-bit
trends) it draws on public disclosures current as of early 2026, cited in the
references. Where a claim is an inference rather than an established fact, the text
says so. History and trend notes are included where they help explain why today's
defaults exist, not as a full chronology of the field.

\keyidea{The goal of an optical systems engineer is not to know every
component. It is to make good engineering decisions under uncertainty using
measurements, physics, and evidence. At gigawatt, multi-generation scale, the
optical interconnect and its lasers are a first-order lever on power, cost, and
reliability. Everything here serves that argument.}

Editorial note: earlier interview-preparation weighting overemphasized
product-readiness processes. The revised structure prioritizes signaling, device
design, packaging, and AI-fabric architecture based on actual interview coverage.
Productization remains as one compact lifecycle chapter; detailed operational and
qualification catalogs are reference content.

Content freeze: prefer cuts, cross-refs, and errata over new frameworks. Structural
relocation and interview-priority rebalancing from the reorientation guide are
allowed and recorded in the errata.
